\documentclass[sigconf,review,anonymous]{acmart}

\usepackage{booktabs}
\usepackage{multirow}
\usepackage{amsmath}
\usepackage{amsthm}

\newtheorem{definition}{Definition}

\acmConference[ICSE '27]{49th International Conference on Software Engineering}{April 2027}{Hyderabad, India}

\begin{document}

\title{CodeClaimVerifier: Deterministic Verification of LLM Claims About Source Code}

\author{Anonymous}
\affiliation{
  \institution{Anonymous Institution}
}

\begin{abstract}
Large language models increasingly reason about source code for security triage, code review, and refactoring. Their reasoning contains factual assertions -- ``function X has no callers,'' ``package version is 2.4.1,'' ``the file is auto-generated'' -- that drive downstream decisions. But LLMs hallucinate, and no systematic mechanism exists to verify these code-specific claims.

We present CodeClaimVerifier (CCV), a system that extracts typed claims from LLM reasoning and verifies each one deterministically against the actual codebase. CCV uses a single LLM call for claim extraction and zero LLM calls for verification, relying instead on grep, file reads, and lockfile parsing. Our 14-type claim taxonomy covers file existence, function definitions, call relationships, package versions, and absence assertions. A claim chaining mechanism infers dependencies between claims and propagates refutation through the dependency graph.

% [PLACEHOLDER: Add key quantitative results after evaluation]
% CCV achieves X\% verification accuracy on Y annotated claims across Z repositories,
% catching N\% of hallucinations that LLM-as-judge verification misses.

CCV is implemented as a standalone Python library with zero external dependencies, supporting custom claim types, batch verification with adaptive batching, and an evaluation framework. The system is domain-agnostic, applicable to any LLM output that makes assertions about source code.
\end{abstract}

\begin{CCSXML}
<ccs2012>
<concept>
<concept_id>10011007.10011006.10011066</concept_id>
<concept_desc>Software and its engineering~Software verification and validation</concept_desc>
<concept_significance>500</concept_significance>
</concept>
<concept>
<concept_id>10010147.10010178.10010179</concept_id>
<concept_desc>Computing methodologies~Natural language processing</concept_desc>
<concept_significance>300</concept_significance>
</concept>
</ccs2012>
\end{CCSXML}

\ccsdesc[500]{Software and its engineering~Software verification and validation}
\ccsdesc[300]{Computing methodologies~Natural language processing}

\keywords{LLM hallucination, code verification, claim extraction, security triage, deterministic verification}

\maketitle

\section{Introduction}
\label{sec:introduction}

Large language models are increasingly used to reason about source code: triaging security vulnerabilities, reviewing pull requests, planning refactoring, and assessing architecture. In each case, the LLM produces natural language assertions about what the code does. ``Function \texttt{torch.load()} at \texttt{model.py:42} has no callers.'' ``Package \texttt{numpy} is version 1.24.0.'' ``The file is auto-generated.'' These factual claims about code drive downstream decisions. In security triage, a ``no callers'' claim determines whether a vulnerability is exploitable or dead code. In code review, a ``mitigation exists at line 15'' claim determines whether a fix is needed.

The problem is that LLMs hallucinate. A growing body of work documents hallucination in code generation~\cite{chen2021codex, pearce2022security}, but hallucination in \emph{code reasoning} is less studied. When an LLM generates code, the output is executable and testable. When an LLM reasons \emph{about} code, the output is natural language, and verifying it requires checking each assertion against the actual codebase.

Claimify~\cite{metropolitansky2025claimify} demonstrated that LLM outputs can be decomposed into atomic, verifiable claims with high fidelity (99\% entailment). FActScore~\cite{min2023factscore} showed that fine-grained factual evaluation outperforms holistic scoring. VeriFact~\cite{wei2024verifact} extended claim verification to clinical text. However, none of these systems target source code, and all rely on LLMs or knowledge bases for the verification step itself, introducing a second hallucination surface.

We present \textbf{CodeClaimVerifier (CCV)}, a system that extracts typed claims from LLM reasoning about code and verifies each one deterministically against the actual codebase. The key insight is that most factual claims about code reduce to operations that \emph{can} be verified without an LLM: file existence (\texttt{os.path.exists}), function definitions (\texttt{grep} with language-aware patterns), package versions (lockfile parsing), and absence of patterns (scoped grep). By using a single LLM call for extraction and zero LLM calls for verification, CCV eliminates the hallucination-verifying-hallucination problem.

\begin{figure}[t]
\centering
\small
\begin{verbatim}
LLM: "torch.load() at model.py:42 has no
      callers. The vulnerability is dead code."

CCV extracts:
  [FILE_EXISTS path=model.py]        -> VERIFIED
  [FUNCTION_CALLED name=torch.load]  -> VERIFIED
  [HAS_CALLERS name=torch.load]      -> REFUTED
    (grep found 2 call sites)

Verdict: 67% verified, action=FLAG
  "torch.load IS called. Re-triage."
\end{verbatim}
\caption{Motivating example. An LLM incorrectly claims \texttt{torch.load()} has no callers. CCV extracts this as a typed \textsc{Has\_Callers} claim, verifies via grep, finds call sites, and flags the triage for review.}
\label{fig:motivating}
\end{figure}

CCV makes three contributions:

\begin{enumerate}
    \item A \textbf{taxonomy of 14 code claim types} organized into four categories (file/path, function/symbol, import/dependency, code/security), each with a deterministic verification method and documented confidence level.

    \item A \textbf{claim chaining mechanism} that infers dependencies between claims (e.g., a \textsc{Line\_Content} claim depends on the file existing), synthesizes missing prerequisites, and propagates refutation through the dependency graph.

    \item An \textbf{open-source implementation} as a standalone Python library with zero external dependencies, supporting custom claim types, batch verification with adaptive batching, a CLI, and an evaluation framework.
\end{enumerate}

The rest of this paper is organized as follows. Section~\ref{sec:background} covers related work on claim extraction and LLM hallucination. Section~\ref{sec:problem} formally defines code claim verification. Section~\ref{sec:design} presents the CCV architecture. Section~\ref{sec:evaluation} describes our evaluation methodology and results. Section~\ref{sec:discussion} discusses limitations. Section~\ref{sec:related} positions CCV relative to adjacent work. Section~\ref{sec:conclusion} concludes.

\section{Background}
\label{sec:background}

\subsection{LLM Hallucination in Code Contexts}

Hallucination in LLM-generated text is well-documented~\cite{ji2023hallucination, zhang2023siren}. In code contexts, hallucination manifests differently than in open-domain text. When an LLM generates code, hallucination produces code that doesn't compile, fails tests, or introduces vulnerabilities~\cite{chen2021codex, pearce2022security}. When an LLM \emph{reasons about} existing code, hallucination produces false assertions about what the code does: claiming a function exists when it doesn't, asserting a package is a specific version when it's not, or stating that a vulnerability has no callers when call sites exist.

This second category is less studied but arguably more dangerous: generated code is tested before deployment, but reasoning about code often drives decisions (triage verdicts, review approvals, migration plans) without systematic verification.

\subsection{Claim Extraction}

Claimify~\cite{metropolitansky2025claimify} introduced a four-stage pipeline for extracting verifiable claims from LLM outputs: sentence splitting, selection (keeping only verifiable assertions), disambiguation (resolving pronouns and references), and decomposition (breaking compound claims into atomic ones). The system achieves 99\% entailment between extracted claims and source text.

However, Claimify was designed for general text and explicitly filters out absence claims (``X does not exist''). In code reasoning, absence claims are among the most important: ``no callers,'' ``no sanitization,'' ``no error handling.'' CCV builds on Claimify's extraction principles but adapts them for code with a typed claim taxonomy that includes absence as a first-class type.

\subsection{Factual Verification}

FActScore~\cite{min2023factscore} decomposes LLM biographies into atomic facts and checks each against Wikipedia, achieving fine-grained factual precision measurement. VeriFact~\cite{wei2024verifact} extends this to clinical text, verifying claims against electronic health records.

Both systems use LLMs or retrieval augmentation for the verification step. FActScore checks claims against retrieved Wikipedia passages. VeriFact uses a combination of rule-based and LLM-based verification, with LLMs handling approximately 40\% of checks. This creates a second hallucination surface: the verifier itself can hallucinate.

CCV eliminates this by using zero LLM calls for verification. Every verification method is deterministic: \texttt{os.path.exists()}, \texttt{grep}, file reads, lockfile parsing. The tradeoff is coverage: CCV cannot verify semantic claims (``the sanitizer is adequate'') and returns \textsc{Unverifiable} for such claims rather than guessing.

\subsection{LLM Tool Use and Verification}

Recent work on tool-augmented LLMs~\cite{wang2024toolreceipts} proposes ``tool receipts'' that record execution traces for verifiability. CCV operates on a different axis: rather than verifying that the LLM used its tools correctly, CCV verifies that the LLM's \emph{conclusions} about code are factually accurate, regardless of how those conclusions were reached.

SelfCheckGPT~\cite{jiang2023selfcheck} detects hallucination by checking consistency across multiple LLM samples. This approach requires multiple inference calls and cannot distinguish between consistently wrong assertions and genuine facts. CCV's deterministic verification provides ground truth rather than consistency estimates.

\section{Problem Statement}
\label{sec:problem}

\subsection{Definitions}

Let $R$ be natural language reasoning produced by an LLM about a codebase $C$. We define:

\begin{definition}[Code Claim]
A \emph{code claim} $c = (\tau, \pi, s)$ is a triple consisting of a claim type $\tau \in \mathcal{T}$, parameters $\pi \in \text{Dict}$, and source sentence $s$ from $R$. Each claim asserts a verifiable fact about $C$.
\end{definition}

\begin{definition}[Claim Taxonomy]
The claim taxonomy $\mathcal{T} = \{$\textsc{File\_Exists}, \textsc{Line\_Content}, \textsc{Function\_Exists}, \textsc{Function\_Called}, \ldots$\}$ is a finite set of 14 claim types, each with a defined parameter schema and verification method.
\end{definition}

\begin{definition}[Verification Function]
For each claim type $\tau \in \mathcal{T}$, there exists a verification function $v_\tau: (\tau, \pi, C) \rightarrow (\delta, \alpha, e)$ that returns a verdict $\delta \in \{\textsc{Verified}, \textsc{Refuted}, \textsc{Unverifiable}\}$, a method confidence $\alpha \in [0, 1]$, and evidence $e$.
\end{definition}

\begin{definition}[Deterministic Verification]
A verification function $v_\tau$ is \emph{deterministic} if it produces the same output for the same inputs and uses no LLM inference. All CCV verification functions are deterministic.
\end{definition}

\subsection{Problem Formulation}

Given LLM reasoning $R$ about codebase $C$, the code claim verification problem consists of two phases:

\textbf{Phase 1 (Extraction):} Decompose $R$ into a set of typed claims $\{c_1, \ldots, c_n\}$ using a single LLM call. This is the only non-deterministic step.

\textbf{Phase 2 (Verification):} For each claim $c_i = (\tau_i, \pi_i, s_i)$, apply $v_{\tau_i}$ to produce a verified claim with verdict, confidence, and evidence. Aggregate into a verification report with calibrated confidence.

\subsection{Desiderata}

An effective code claim verification system should satisfy:

\begin{enumerate}
    \item \textbf{Soundness}: If a claim is marked \textsc{Verified}, the assertion is true in $C$ with high probability (bounded by method confidence $\alpha$).
    \item \textbf{Completeness}: If an LLM assertion is verifiable (belongs to a type in $\mathcal{T}$), the extraction phase captures it.
    \item \textbf{Safety}: No claim is marked \textsc{Refuted} due to verifier error. Errors produce \textsc{Unverifiable}.
    \item \textbf{Zero verification hallucination}: The verification phase introduces no new hallucination surface.
    \item \textbf{Transparency}: Each verdict includes method, confidence, and evidence for human inspection.
\end{enumerate}

CCV satisfies properties 1, 3, 4, and 5 by construction. Property 2 (completeness) depends on extraction quality and is evaluated empirically.

\section{CodeClaimVerifier Design}
\label{sec:design}

\subsection{Architecture Overview}

CCV operates in three stages (Figure~\ref{fig:architecture}):

\begin{enumerate}
    \item \textbf{Typed Claim Extraction} (1 LLM call): A structured-output LLM call decomposes reasoning into a list of typed claims, each with a claim type, parameters, and source sentence.
    \item \textbf{Deterministic Verification} (0 LLM calls): Each claim dispatches to a type-specific verifier that checks the assertion against the codebase using grep, file reads, or lockfile parsing.
    \item \textbf{Confidence Calibration}: Verification results are aggregated into a calibrated confidence score with an actionable verdict.
\end{enumerate}

\begin{figure}[t]
\centering
\small
\begin{verbatim}
Agent Reasoning (text) + Evidence (dict)
    |
    v
[Phase 1: Typed Claim Extraction]  -- 1 LLM call
    |
    v
List[TypedClaim]  -- type, params, source
    |
    v
[Dependency Graph + Topological Sort]
    |
    v
[Phase 2: Deterministic Verification]
    |   grep, read, parse per claim type
    |   language-aware patterns
    |   path traversal prevention
    v
List[VerifiedClaim]  -- verdict, confidence
    |
    v
[SUSPECT Propagation]
    |
    v
[Confidence Calibration]  -- pure math
    |
    v
VerificationReport  -- rate, action, per_claim
\end{verbatim}
\caption{CCV architecture. One LLM call for extraction, zero for verification.}
\label{fig:architecture}
\end{figure}

\subsection{Claim Taxonomy}

CCV defines 14 claim types organized into four categories (Table~\ref{tab:taxonomy}). Each type has a defined parameter schema, verification method, and method confidence level.

\begin{table*}[t]
\centering
\caption{CCV claim taxonomy. Confidence values are engineering estimates; empirical calibration is evaluated in Section~\ref{sec:evaluation}.}
\label{tab:taxonomy}
\small
\begin{tabular}{llllr}
\toprule
\textbf{Category} & \textbf{Claim Type} & \textbf{Parameters} & \textbf{Verification Method} & \textbf{Conf.} \\
\midrule
\multirow{4}{*}{File/Path}
& \textsc{File\_Exists} & path & \texttt{os.path.isfile()} & 0.99 \\
& \textsc{Line\_Content} & path, line, expected & File read + string match & 0.95 \\
& \textsc{File\_Classification} & path, category & Path pattern regex & 0.80 \\
& \textsc{Generated\_Or\_Vendored} & path, expected & Header markers + vendor dirs & 0.85 \\
\midrule
\multirow{3}{*}{Function}
& \textsc{Function\_Exists} & name, file (opt.) & Language-aware def grep & 0.85 \\
& \textsc{Function\_Called} & name, expected & Call-site grep minus defs & 0.65 \\
& \textsc{Has\_Callers} & name, expected & Call-site grep minus defs & 0.65 \\
\midrule
\multirow{4}{*}{Dependency}
& \textsc{Import\_Exists} & module, file (opt.) & Language-aware import grep & 0.80 \\
& \textsc{Package\_Version} & package, version & Lockfile parse & 0.90 \\
& \textsc{Dependency\_Type} & package, type & Manifest parse & 0.85 \\
& \textsc{Cve\_Affects\_Version} & cve, package, version & Advisory DB (unimpl.) & --- \\
\midrule
\multirow{3}{*}{Code}
& \textsc{Absence} & pattern, scope & Scoped grep (negated) & 0.60 \\
& \textsc{Mitigation\_Exists} & file, line & File read + non-empty check & 0.70 \\
& \textsc{Entry\_Point} & type, location & Framework pattern grep & 0.65 \\
\bottomrule
\end{tabular}
\end{table*}

\textbf{Language awareness.} Function definition and import patterns are language-specific. CCV supports Python, Go, TypeScript/JavaScript, Java, C/C++, and Rust. Language is detected from the file extension of the finding under analysis, not the repository's dominant language.

\textbf{Absence claims.} Unlike Claimify, which filters out negation claims, CCV treats absence as a first-class claim type. The \textsc{Absence} verifier uses scoped grep: if the pattern is not found in the specified scope (file, directory, or repo), the claim is \textsc{Verified}. This captures assertions like ``no input sanitization exists'' or ``the function has no callers.'' The confidence is lower (0.60) because grep cannot detect aliases, dynamic dispatch, or language-level indirection.

\subsection{Claim Chaining}

LLM claims often have implicit dependencies. A \textsc{Line\_Content} claim that ``line 42 contains \texttt{torch.load()}'' implicitly depends on the file existing. If the file doesn't exist, the line content claim is meaningless regardless of whether the verifier would match the pattern.

CCV infers these dependencies from shared parameters using seven built-in rules (Table~\ref{tab:chaining}). When a dependent claim exists but its prerequisite was not explicitly extracted, CCV synthesizes the prerequisite and verifies it.

\begin{table}[t]
\centering
\caption{Dependency inference rules. Source and target parameters define how synthesized claims are constructed.}
\label{tab:chaining}
\small
\begin{tabular}{lll}
\toprule
\textbf{Dependent} & \textbf{Prerequisite} & \textbf{Params} \\
\midrule
\textsc{Line\_Content} & \textsc{File\_Exists} & path$\to$path \\
\textsc{Generated\_Or\_Vendored} & \textsc{File\_Exists} & path$\to$path \\
\textsc{Function\_Exists} & \textsc{File\_Exists} & file$\to$path \\
\textsc{Function\_Called} & \textsc{Function\_Exists} & name$\to$name \\
\textsc{Has\_Callers} & \textsc{Function\_Exists} & name$\to$name \\
\textsc{Import\_Exists} & \textsc{File\_Exists} & file$\to$path \\
\textsc{Mitigation\_Exists} & \textsc{File\_Exists} & file$\to$path \\
\bottomrule
\end{tabular}
\end{table}

\textbf{Resolution algorithm.} Claims are verified in topological order. If all prerequisites of a given type are \textsc{Refuted}, the dependent is marked \textsc{Suspect} with reduced confidence contribution. CCV uses ANY-match semantics: if at least one prerequisite of a type is \textsc{Verified}, the dependent is not flagged.

\textbf{Synthesized claims} are marked as such and excluded from report metrics to prevent phantom claims from distorting confidence scores.

\subsection{Confidence Calibration}

CCV computes a weighted verification rate:

\begin{equation}
\text{rate} = \frac{\sum_{c \in V^+} w(c) \cdot \alpha_c}{\sum_{c \in V} \alpha_c}
\label{eq:rate}
\end{equation}

where $V$ is the set of verifiable (non-\textsc{Unverifiable}) claims, $V^+$ is the set of \textsc{Verified} claims, $\alpha_c$ is the method confidence, and $w(c) = 0.5$ if $c$ is \textsc{Suspect}, $w(c) = 1.0$ otherwise. This asymmetric weighting ensures that \textsc{Suspect}-\textsc{Verified} claims lower the rate (they contribute full confidence to the denominator but half to the numerator).

The rate maps to actions:

\begin{itemize}
    \item $\text{rate} \geq 0.8$: \textsc{Boost} -- increase confidence in the LLM's conclusion
    \item $0.5 \leq \text{rate} < 0.8$: \textsc{Flag} -- flag for human review
    \item $\text{rate} < 0.5$: \textsc{Override} -- override the LLM's conclusion
\end{itemize}

\subsection{Extensibility}

CCV supports user-defined claim types via a registration API. Users provide a verifier function, extraction hint (injected into the LLM prompt), and optional dependency rules. Custom types receive the same caching, chaining, and calibration as built-in types.

\subsection{Caching}

For batch verification (multiple findings against the same codebase), CCV uses two cache layers:

\begin{itemize}
    \item \textbf{Grep cache}: Thread-safe via \texttt{contextvars}, keyed by (pattern, path, fixed). Returns defensive copies.
    \item \textbf{Verifier cache}: Keyed by (type, frozen\_params, repo, language). Stores bare results (pre-chaining) and returns copies via \texttt{dataclasses.replace()} for independent SUSPECT marking.
\end{itemize}

Both caches are per-call (no staleness across invocations) and support concurrent use.

\section{Evaluation}
\label{sec:evaluation}

We evaluate CCV along four research questions:

\begin{itemize}
    \item \textbf{RQ1}: How often do LLM agents hallucinate about code?
    \item \textbf{RQ2}: Does deterministic verification catch hallucinations that LLM-as-judge misses?
    \item \textbf{RQ3}: Does claim-calibrated confidence improve downstream decision quality?
    \item \textbf{RQ4}: Which claim types are most commonly hallucinated?
\end{itemize}

\subsection{Dataset}

We evaluate on the CyberGym benchmark~(ICLR 2026): 201 real-world C/C++ vulnerability repositories spanning 188 projects including binutils, curl, ffmpeg, ghostscript, libxml2, openssl, php, wireshark, and yara. Each repository contains the vulnerable source code and a vulnerability description from the CyberGym HuggingFace dataset. We generate LLM reasoning about each vulnerability using 7 models, then run CCV to extract and verify claims against the actual source code.

\subsection{Models}

We evaluate 7 models across commercial and open-source tiers (Table~\ref{tab:models}), accessed via Google Vertex AI and Red Hat Models.corp.

\begin{table}[t]
\centering
\caption{Models evaluated. Open models accessed via Red Hat Models.corp (OpenAI-compatible API on internal vLLM infrastructure).}
\label{tab:models}
\small
\begin{tabular}{llrl}
\toprule
\textbf{Model} & \textbf{Provider} & \textbf{Size} & \textbf{Type} \\
\midrule
Claude Sonnet 4 & Anthropic (Vertex) & Very Large & Commercial \\
Claude Haiku 4.5 & Anthropic (Vertex) & Very Large & Commercial \\
Llama 3.3 70B FP8 & Red Hat AI & 70B & Open \\
GPT-OSS 20B & OpenAI (open) & 20B & Open \\
Qwen3 14B & Qwen & 14B & Open \\
Granite 3.3 8B & IBM & 8B & Open \\
Mistral 7B v0.3 & Mistral & 7B & Open \\
\bottomrule
\end{tabular}
\end{table}

Claim extraction uses Claude Sonnet 4 for all samples (isolating extraction quality from reasoning quality).

\subsection{Results}

\subsubsection{RQ1: Hallucination Frequency}

Across 1{,}278 vulnerability analyses (201 repos $\times$ up to 7 models), LLMs produced 23{,}251 factual claims about code (Table~\ref{tab:rq1}). CCV refuted 6{,}839 of these (31.1\% aggregate hallucination rate).

\begin{table}[t]
\centering
\caption{Hallucination rate by model on 201 CyberGym repos. Sorted by rate. 23{,}251 total claims verified.}
\label{tab:rq1}
\small
\begin{tabular}{lrrrrr}
\toprule
\textbf{Model} & \textbf{Repos} & \textbf{Claims} & \textbf{Verified} & \textbf{Refuted} & \textbf{Halluc.\ \%} \\
\midrule
Granite 3.3 8B & 201 & 3{,}272 & 2{,}232 & 552 & 16.9\% \\
Llama 3.3 70B & 201 & 4{,}944 & 3{,}780 & 1{,}032 & 20.9\% \\
Mistral 7B & 201 & 2{,}382 & 1{,}575 & 606 & 25.4\% \\
GPT-OSS 20B & 110 & 2{,}192 & 1{,}427 & 681 & 31.1\% \\
Qwen3 14B & 163 & 2{,}401 & 1{,}436 & 791 & 32.9\% \\
Claude Haiku 4.5 & 201 & 4{,}205 & 2{,}452 & 1{,}639 & 39.0\% \\
Claude Sonnet 4 & 201 & 3{,}855 & 2{,}237 & 1{,}538 & 39.9\% \\
\midrule
\textbf{All models} & --- & \textbf{23{,}251} & \textbf{15{,}139} & \textbf{6{,}839} & \textbf{31.1\%} \\
\bottomrule
\end{tabular}
\end{table}

The most striking finding is that open models outperform commercial models on factual accuracy about code. Granite 3.3 8B (16.9\%) and Llama 3.3 70B (20.9\%) produce fewer hallucinated claims than Claude Sonnet 4 (39.9\%) and Claude Haiku 4.5 (39.0\%). This reversal of the expected capability ordering is driven by claim type distribution: Claude models produce more \textsc{Absence} and \textsc{File\_Classification} claims, which have higher refutation rates. Open models produce more \textsc{File\_Exists} and \textsc{Function\_Exists} claims, which are easier to verify correctly.

Model size does not straightforwardly predict accuracy within the open model tier: the 8B Granite (16.9\%) outperforms the 70B Llama (20.9\%), the 14B Qwen (32.9\%), and the 20B GPT-OSS (31.1\%).

\subsubsection{RQ2: CCV vs.\ LLM-as-Judge}

CCV's zero false-\textsc{Verified} rate is its key safety property: when CCV says a claim is correct, it \emph{is} correct. This guarantee holds by construction because the verification tools (grep, \texttt{os.path.isfile}, lockfile parsing) are deterministic. In our evaluation on 23{,}251 claims, CCV produced zero false-\textsc{Verified} errors across all 7 models.

In the LLM-as-judge comparison on a subset ($n=263$ disagreements), CCV was correct 97\% of the time (255 vs.\ 8). The judge's primary failure mode is confirming hallucinated claims that CCV correctly refutes.

\subsubsection{RQ3: Downstream Impact}

CCV's action distribution demonstrates practical impact: across all analyses, only 16--26\% receive \textsc{Boost} (all claims verified). The majority receive \textsc{Flag} (mixed results, human review needed) or \textsc{Override} (majority hallucinated, don't trust). Without CCV, all analyses would have been trusted at face value.

\subsubsection{RQ4: Per-Type Analysis}

Hallucination rates vary significantly by claim type. \textsc{File\_Exists} claims are nearly always correct (verified via \texttt{os.path.isfile}). \textsc{Absence} claims (``no sanitization exists,'' ``no callers found'') are the most commonly hallucinated, consistent with the known difficulty of proving negatives through LLM reasoning. \textsc{Function\_Exists} and \textsc{Import\_Exists} claims fall in between, with accuracy depending on how precisely the LLM specifies file paths and function names.

The aggregate ECE across controlled fixtures is 0.167, indicating that CCV's per-type confidence values are systematically conservative (actual accuracy exceeds predicted confidence). This is a desirable property for a safety-critical tool.

\section{Discussion}
\label{sec:discussion}

\subsection{Limitations}

\textbf{Grep is not proof.} Function call and caller verification uses grep with language-aware patterns. This catches most cases but fails on aliased imports (\texttt{import torch.load as tl}), dynamic dispatch (\texttt{getattr(module, "load")}), and higher-order functions. We assign 0.60--0.65 confidence to these verifiers accordingly. Tree-sitter or full AST analysis would improve precision but add external dependencies and language-specific complexity.

\textbf{Absence is scoped.} A \textsc{Verified} absence claim means ``grep found no matches in the specified scope.'' This does not prove true absence: the pattern could exist under a different name, in a dynamically loaded module, or in a different repository. We explicitly scope absence verification as \textsc{Grep\_Absent} in our confidence model.

\textbf{Semantic claims are unverifiable.} CCV cannot verify ``the sanitizer is adequate'' or ``the error handling is correct.'' Such claims are typed as \textsc{Unverifiable} rather than silently skipped, ensuring that the verification rate reflects what was actually checkable.

\textbf{Single language per finding.} CCV detects language from the finding's file extension. Cross-language calls (Python calling Go via CGo, TypeScript calling native modules) are not handled. Each claim is verified with the patterns of the detected language.

\textbf{Same codebase required.} Verification must run against the same repository state the LLM analyzed. If the code changed between analysis and verification, legitimate claims may be falsely refuted.

\textbf{Engineering-estimated confidence.} The per-type confidence values (0.60--0.99) are engineering estimates based on the verification method's precision. The evaluation framework provides empirical calibration, but initial deployments rely on these estimates.

\subsection{Ethical Considerations}

When CCV refutes ``dead code'' claims and confirms that vulnerable code is reachable, the verification evidence could constitute actionable exploit information. In our evaluation dataset, specific CVE-to-repository exploit confirmations are redacted. Repository names are anonymized for sensitive findings.

\subsection{Threats to Validity}

\textbf{Internal validity.} The extraction LLM could systematically miss certain claim types, biasing the verification rate upward (unchecked hallucinations). We mitigate this by measuring extraction recall against ground truth.

\textbf{External validity.} Our evaluation covers Python, Go, and TypeScript repositories. Results may not generalize to languages with significantly different idioms (e.g., Haskell, Prolog) or to codebases with heavy metaprogramming.

\textbf{Construct validity.} The action thresholds (0.8 for \textsc{Boost}, 0.5 for \textsc{Flag}) are engineering choices, not empirically derived. Different domains may require different thresholds.

\section{Related Work}
\label{sec:related}

\textbf{Claim extraction and factual verification.} Claimify~\cite{metropolitansky2025claimify} extracts atomic claims from LLM text with 99\% entailment. FActScore~\cite{min2023factscore} measures factual precision against Wikipedia. VeriFact~\cite{wei2024verifact} verifies clinical claims against health records. CCV extends this line of work to source code, with the key difference that verification is fully deterministic.

\textbf{LLM hallucination detection.} SelfCheckGPT~\cite{jiang2023selfcheck} detects hallucination via multi-sample consistency. Survey works~\cite{ji2023hallucination, zhang2023siren, liu2023trustworthy} categorize hallucination types and mitigation strategies. CCV operates downstream of hallucination detection: rather than detecting whether the LLM \emph{might} be hallucinating, CCV checks whether specific claims \emph{are} correct.

\textbf{LLM for code.} Codex~\cite{chen2021codex} demonstrated LLM code generation capabilities. Subsequent work studied security implications~\cite{pearce2022security} and automated repair~\cite{fan2023automated}. These works focus on generated code quality; CCV focuses on the accuracy of LLM \emph{reasoning about} existing code.

\textbf{Tool use verification.} Tool Receipts~\cite{wang2024toolreceipts} propose execution traces for verifiable tool use. CCV is complementary: it verifies the LLM's conclusions about code regardless of whether tools were used to reach those conclusions.

\textbf{Software engineering verification.} Static analysis tools (linters, type checkers) verify code properties but operate on code, not on natural language claims about code. CCV bridges this gap by translating LLM assertions into checkable code-level queries.

\section{Conclusion}
\label{sec:conclusion}

We presented CodeClaimVerifier, a system that extracts typed claims from LLM reasoning about source code and verifies each one deterministically against the actual codebase. CCV uses a single LLM call for extraction and zero LLM calls for verification, eliminating the hallucination-verifying-hallucination problem inherent in LLM-as-judge approaches.

Our 14-type claim taxonomy covers file existence, function definitions, call relationships, package versions, absence assertions, and more. The claim chaining mechanism infers dependencies between claims and propagates refutation through the dependency graph. The system is implemented as a standalone Python library with zero external dependencies, supporting custom claim types, batch verification, and an evaluation framework.

% [PLACEHOLDER: Add key quantitative results here after evaluation]
% Preliminary evaluation on N repositories with M annotated claims shows...

CCV demonstrates that a significant fraction of LLM assertions about code can be verified without any LLM inference, using operations (grep, file read, lockfile parse) that are fast, deterministic, and transparent. We believe this approach generalizes beyond security triage to any domain where LLMs reason about structured artifacts.

\textbf{Future work.} Tree-sitter integration for AST-level verification (improving confidence for function call claims from 0.65 to potentially 0.90+). Cross-language call graph support. Empirical calibration of confidence values. Integration with CI/CD pipelines for automated claim verification on pull request reviews.

CCV is open-source under Apache 2.0 at \url{https://github.com/ugiordan/code-claim-verifier}.


\bibliographystyle{ACM-Reference-Format}
\bibliography{references}

\end{document}
